% ECONOMICS_PROBLEM: {"id": "C73-113", "title": "Constant regret in general-sum Markov games", "jel_code": "C73", "source_id": "OP-0151"}
% FIRST_POSTED: 2026-10-11
% ATTEMPT_STATUS: solved
% ENTRY_KIND: research_attempt
\documentclass[11pt]{article}
\usepackage[a4paper,margin=25mm]{geometry}
\usepackage[T1]{fontenc}
\usepackage{lmodern,microtype,amsmath,amssymb,amsthm}
\usepackage[hidelinks]{hyperref}
\newtheorem{theorem}{Theorem}
\newtheorem{lemma}{Lemma}
\title{Constant regret with finite preparation in Markov games}
\author{Ji Ho Bae}
\date{October 2026}
\hypersetup{pdftitle={Constant regret with finite preparation in Markov games},pdfauthor={Ji Ho Bae}}
\begin{document}
\maketitle

\paragraph{Authorship and provenance.}
Author: Ji Ho Bae. Prepared with GPT Codex assistance; the precise serving
revision was not exposed. The complete original argument is retained in
the \href{https://github.com/jbaelaw/markov-game-constant-regret-proof/blob/bae26db857293585b886da154ab9e4b2fcc16949/paper/Proof.tex}{immutable TeX source} and
\href{https://github.com/jbaelaw/markov-game-constant-regret-proof/blob/bae26db857293585b886da154ab9e4b2fcc16949/paper/Proof.pdf}{proof PDF}. This submission adds catalogue metadata,
provenance and a completion statement to that argument. Its preparation
included exact-target and primary-source checks and separate mathematical
reviews. This is a complete proof claim submitted for independent expert
review; community acceptance remains pending.

\begin{abstract}
We answer the computational formulation of catalogue problem C73-113
under its rational input convention and observable-action feedback.
Players reconstruct their local tables, communicate through observed
actions, and compute a rational coarse equilibrium over pure Markov
plans. A maximum-deficit schedule tracks its weights with uniformly
bounded error. This gives constant cumulative external regret against
all history-dependent within-episode policies, with polynomial work per
episode. The finite preparation period is included in the regret bound
and may be very large as a function of the game.
\end{abstract}

\section{The stated target}
There are $n$ players, a finite state set $S$, nonempty finite action
sets $A_i$, a horizon $H\geq1$, and a fixed initial state $s_1$.
Rewards $r_{i,h}(s,a)$ lie in $[0,1]$, and $P_h(\cdot\mid s,a)$ is the
next-state distribution. We use the catalogue's computational convention:
the finite game data are rational and binary encoded
\cite[Common conventions, Computational representations]{catalogue}.
Let $L$ be their description length. The labeled finite structure is
public, while reward and transition entries are initially accessible
through the prescribed feedback only.

At a queried $(h,s,a_i)$, player $i$ receives its own reward averaged
against the opponents' current Markov policies and the corresponding
averaged transition distribution, both exactly. All realized states and actions are observed. A behavioral
policy may depend on the entire observed history within an episode.
Let $\Pi_i$ be this policy class, and let $V_i(\pi)$ be expected total
reward in one episode. Thus $0\leq V_i(\pi)\leq H$. The target is
\[
 R_i(T)=\sup_{\tau_i\in\Pi_i}\sum_{t=1}^T
       \bigl(V_i(\tau_i,\pi^t_{-i})-V_i(\pi^t)\bigr)\leq C_G
 \qquad(T\geq1),
\]
where $\tau_i$ is one fixed within-episode policy and $C_G$ is finite
and independent of $T$. Neither $C_G$ nor total preparation time is
required to be polynomial in $L$. The resource requirement is polynomial
work and polynomially many feedback queries in each episode, as a
function of the game description and the current episode number.

The source motivates constant regret and supplies the own-feedback
oracles \cite[\S\S2.3--2.5 and 5]{yb}. The catalogue specifies the
ordinary external-regret target above. The distinction matters: the
construction below uses a finite initial calculation and action signals,
and its bound is stated explicitly in terms of that calculation.

\begin{theorem}\label{main}
One deterministic decentralized procedure, independent of the final
number of episodes, achieves the target for every game in this model.
Every played profile is pure Markov. Each player makes at most one
feedback query per episode, and its bit operations are bounded by a
fixed polynomial in $L$ and the current episode number.
More precisely, the procedure has a finite preparation duration $T_0$
and computes a distribution with finite support size $K$, for which
\[
                  R_i(T)\leq H T_0+2H(K-1)
                  \qquad\text{for every }i\text{ and }T\geq1.
\]
\end{theorem}

If every action set is a singleton, play the unique profile forever;
all regrets are zero. Otherwise, call players with $|A_i|\geq2$ active.
A remaining player's regret is always zero, so its reward table is
irrelevant to the computation below.

\section{A rational distribution covering every deviation}
The played plans depend only on $(h,s)$; comparison policies may also
depend on earlier states and actions.
Let $\Sigma_i$ be all pure Markov plans $[H]\times S\to A_i$, and put
$\Sigma=\prod_i\Sigma_i$. Let $D_i$ be all pure policies on the finite
set of player $i$'s pre-action histories, including histories of zero
probability. Histories contain past states and joint actions and the
current state, before the current actions are selected. Define
\[
 g_{i,d}(\sigma)=V_i(d,\sigma_{-i})-V_i(\sigma)
 \qquad(d\in D_i,\ \sigma\in\Sigma).
\]
These are rational numbers computable by finite enumeration of play
histories. Consider the finite linear system
\begin{equation}\label{lp}
 \mu_\sigma\geq0,\qquad \sum_{\sigma\in\Sigma}\mu_\sigma=1,
 \qquad \sum_{\sigma\in\Sigma}\mu_\sigma g_{i,d}(\sigma)\leq0
 \quad(i\text{ active},\ d\in D_i).
\end{equation}

This system is nonempty. To prove it, set $v_{i,H+1}(s)=0$ and work
backward. At $(h,s)$ choose a mixed Nash equilibrium of the finite
stage game with payoffs
\[
 r_{i,h}(s,a)+\sum_{s'}P_h(s'\mid s,a)v_{i,h+1}(s'),
\]
and let $v_{i,h}(s)$ be its expected payoff. Such equilibria exist by
Nash's finite-game theorem \cite{nash}. The resulting behavioral Markov
profile $x$ resists every history-dependent unilateral deviation.
Indeed, at any history ending in $(h,s)$, the opponents' current actions
have the prescribed stage distribution. By induction, a deviator's
future payoff is at most $v_{i,h+1}$ at the next state. The stage Nash
inequality then bounds its current reward plus continuation by
$v_{i,h}(s)$. This proves the assertion backward through all histories.

Independently sample an action from $x_{i,h}(s)$ for each player and
each $(h,s)$, thereby obtaining a random member of $\Sigma$.
Its law implements $x$, including against any fixed history-dependent
deviation: a given stage-state pair is visited at most once, and its
unused action sample is independent of the samples that determined
the past. The induced distribution on $\Sigma$ satisfies
\eqref{lp}. It is used only to prove nonemptiness.

The feasible set is a nonempty compact rational polytope, so it has a
rational vertex. For completeness, successively maximizing each
coordinate on the preceding maximizing face produces a vertex.
At a vertex, active constraint normals span the whole coordinate space;
otherwise a small perturbation in both directions orthogonal to them
would preserve all constraints and contradict extremality. A full-rank
active rational system therefore determines that vertex by Gaussian
elimination. This also gives a finite exact algorithm: enumerate all
full-rank active systems, solve them, test all constraints, and keep the
first passing solution in a fixed public enumeration order.

Discard zero weights and write its support as
$\sigma^1,\ldots,\sigma^K$, with positive rational weights
$\mu_1,\ldots,\mu_K$. For any behavioral $\tau_i$, independently
pre-sampling its action at each finite history gives a distribution on
$D_i$ that implements $\tau_i$ against every fixed $\sigma_{-i}$.
The distribution is the same for all $\sigma$. Averaging the corresponding
rows of \eqref{lp} proves
\begin{equation}\label{allpolicies}
 \sum_{j=1}^K\mu_j
       \bigl(V_i(\tau_i,\sigma^j_{-i})-V_i(\sigma^j)\bigr)\leq0
       \qquad(\tau_i\in\Pi_i).
\end{equation}
Singleton-action players satisfy this with equality. Thus the supported
profiles are Markov, while the comparison class retains all histories.

\section{Legal decentralized preparation}
Fix public orders on players, states, actions and table entries, and
one baseline action for each player. Enumerate $(h,s,a)\in[H]\times
S\times\prod_iA_i$, reserving one episode for each triple.
In that episode every player $j$ uses the constant Markov plan $a_j$.
Each active player $i$ queries $(h,s,a_i)$ once. The opponents' policy
is a point mass, so the reply is exactly $r_{i,h}(s,a)$ and
$P_h(\cdot\mid s,a)$. After these finitely many episodes, each active
player has its own reward table and the same full transition table.
All queries concern the actual current opponent policy, including when
the queried state is unreachable in that episode.

Active players next broadcast their reward tables, in public player
order. For a sender $i$, fix two distinct actions labeled zero and one.
It sends a bit through its action at $(1,s_1)$; every other action follows
the baseline. Everyone observes the bit with certainty, independently
of subsequent state transitions. These signaling plans are pure Markov.
No player queries another player's reward oracle.

Use canonical finite binary encodings of rational entries and the fixed
table order. A payload of length $b\geq1$ is prefixed by
$1^\ell0\,\operatorname{bin}(b)$, where
$\ell=\lfloor\log_2 b\rfloor+1$ and $\operatorname{bin}(b)$ has $\ell$
bits. Receivers count the initial ones, read the length and then the
payload. All players therefore know exactly when the next sender starts.
All these strings have finite length, polynomial in $L$ under the stated
finite-table encoding. Preparing, sending and reading them has
polynomial cost per episode. A sole active player already has every
relevant reward table; the same protocol is harmless in that case.

Every active player now has identical input for \eqref{lp}. Each runs
the same deterministic exact program just described, executing one
transition of a fixed binary Turing machine per episode while playing
the baseline plan. This finite program eventually halts, in the same
episode for every active player. Its simulation includes generating the
linear system, solving it and writing the support plans and weights.
Let $T_0$ count all exploration, signaling and computation episodes,
including input and output preparation. It is finite and depends only
on the game and the fixed procedure.

\section{A schedule with uniformly bounded error}
For $m\geq0$, let $N_j(m)$ count uses of $\sigma^j$ in the first $m$
episodes after preparation, and put
\[
                     d_j(m)=m\mu_j-N_j(m).
\]
Start with all counts zero. In episode $m\geq1$, choose the least index
maximizing $d_j(m-1)+\mu_j$, play $\sigma^j$, and increment its count.
All active players have the same weights, counts and tie rule, so they
select the same joint plan and execute their own components.

For every $m$, the deficits sum to zero and each is greater than $-1$.
The latter holds initially. Before the next selection, the numbers
$d_j(m-1)+\mu_j$ sum to one, so the selected number is positive.
Subtracting one leaves it greater than $-1$; unselected numbers have
only increased. Induction proves the claim. If $K=1$ the deficit is zero.
Otherwise at most $K-1$ coordinates are negative, and each has magnitude
less than one. Zero total sum consequently gives
\begin{equation}\label{discrepancy}
                       \sum_j|d_j(m)|\leq2(K-1).
\end{equation}

\newpage
Fix any player $i$ and any $\tau_i\in\Pi_i$, and abbreviate
$g_j=V_i(\tau_i,\sigma^j_{-i})-V_i(\sigma^j)$.
We have $|g_j|\leq H$ and $\sum_j\mu_jg_j\leq0$ by
\eqref{allpolicies}. Its total deviation gain after preparation is
\[
 \sum_jN_j(m)g_j
   =m\sum_j\mu_jg_j-\sum_jd_j(m)g_j
   \leq H\sum_j|d_j(m)|\leq2H(K-1).
\]
Each preparation episode contributes at most $H$. This proves the bound
in Theorem~\ref{main}, also when $T\leq T_0$. The bound is uniform in
$\tau_i$, so taking its supremum proves the claim. By linearity, the
ordinary uniform empirical mixture has coarse-deviation gain at most
$\{H T_0+2H(K-1)\}/T$.

\section{Per-episode resources}
Exploration uses one own-feedback query per active player per episode;
later phases use none. Communication and its counters have polynomial
bit cost. During the simulated exact computation, at most
$O(\operatorname{poly}(L)+t)$ tape cells have been initialized by episode
$t$. Even direct tape scans therefore implement each simulated transition
with polynomially many bit operations in $L,t$.

The final support and weights were written before episode $T_0$.
Their total encoded length, and hence $K$, are bounded by the same tape
bound. Write $\mu_j=p_j/q_j$ with $q_j>0$. The next comparison value is
\[
          d_j(m-1)+\mu_j
             =\frac{m p_j-N_j(m-1)q_j}{q_j}.
\]
Counts use $O(\log m)$ bits. Exact cross multiplication, scanning the
stored support and retrieving the selected local Markov plan all have
polynomial cost in its stored length and $\log m$, hence in $L,t$.
This establishes a single per-episode polynomial bound for the entire
procedure. All preparation is charged to $T_0$; the resulting
$C_G=H T_0+2H(K-1)$ is finite but may grow rapidly with the game size.
The algorithm never uses the final episode count $T$.

\section*{Completion Estimate}
100\%
The stated result resolves the original catalogue question under its full
model and quantifiers. This is the author's complete proof claim, pending
independent expert review.

\begin{thebibliography}{9}
\bibitem{catalogue} Economics problem C73-113, Constant regret in
general-sum Markov games. Statement and Common conventions,
Computational representations. Catalogue snapshot at commit
\texttt{3cf5ae3}.
\href{https://github.com/mehmetmars7/Erdosproblems-llm-hunter/blob/3cf5ae39f671cb748fa7ba6d0ccc253611620f9d/attacks/open_problems/economics/statements/C73-113.tex}{Immutable problem statement}.
\bibitem{yb} A. E. Yorulmaz and T. Ba\c{s}ar.
Near Optimal Convergence to Coarse Correlated Equilibrium in
General-Sum Markov Games. 2025.
\href{https://arxiv.org/html/2511.02157v1}{arXiv:2511.02157v1},
\S\S2.3--2.5 and 5.
\bibitem{nash} J. F. Nash, Jr. Equilibrium points in n-person games.
\emph{Proceedings of the National Academy of Sciences} 36(1)
(1950), 48--49.
\href{https://doi.org/10.1073/pnas.36.1.48}{doi:10.1073/pnas.36.1.48}.
\end{thebibliography}
\end{document}
